% !tp-example Basic usage
\documentclass{article}
\begin{document}
\wcTpSingleChoiceQuestion[answer=1,random]{
  \begin{description}
    \item[{Title}] Geography
    \item[{Prompt}] What is the capital of France?
    \item[{Form}] \begin{itemize}
      \item Paris
      \item London
      \item Berlin
    \end{itemize}
    \item[{Feedback}] \begin{itemize}
      \item Correct.
      \item London is in the UK.
      \item Berlin is in Germany.
    \end{itemize}
    \item[{Solution}] Paris
  \end{description}
}
\end{document}

% !tp-example Javascript API
\documentclass{article}
\begin{document}
\wcScript[type=module]|
import "/components/single-choice-question/single-choice-question.js";

const example = document.createDocumentFragment();

const element1 = document.createElement("tp-single-choice-question");
element1.setAttribute("answer", "1");
element1.setAttribute("random", "");
element1.append(document.createTextNode("\n    "));

const element2 = document.createElement("dl");
element2.append(document.createTextNode("\n      "));

const element3 = document.createElement("dt");
element3.append(document.createTextNode("Title"));
element2.append(element3);

const element4 = document.createElement("dd");
element4.append(document.createTextNode("Geography"));
element2.append(element4);
element2.append(document.createTextNode("\n      "));

const element5 = document.createElement("dt");
element5.append(document.createTextNode("Prompt"));
element2.append(element5);

const element6 = document.createElement("dd");
element6.append(document.createTextNode("What is the capital of France?"));
element2.append(element6);
element2.append(document.createTextNode("\n      "));

const element7 = document.createElement("dt");
element7.append(document.createTextNode("Form"));
element2.append(element7);

const element8 = document.createElement("dd");

const element9 = document.createElement("ul");

const element10 = document.createElement("li");
element10.append(document.createTextNode("Paris"));
element9.append(element10);

const element11 = document.createElement("li");
element11.append(document.createTextNode("London"));
element9.append(element11);

const element12 = document.createElement("li");
element12.append(document.createTextNode("Berlin"));
element9.append(element12);
element8.append(element9);
element2.append(element8);
element2.append(document.createTextNode("\n      "));

const element13 = document.createElement("dt");
element13.append(document.createTextNode("Feedback"));
element2.append(element13);

const element14 = document.createElement("dd");

const element15 = document.createElement("ul");

const element16 = document.createElement("li");
element16.append(document.createTextNode("Correct."));
element15.append(element16);

const element17 = document.createElement("li");
element17.append(document.createTextNode("London is in the UK."));
element15.append(element17);

const element18 = document.createElement("li");
element18.append(document.createTextNode("Berlin is in Germany."));
element15.append(element18);
element14.append(element15);
element2.append(element14);
element2.append(document.createTextNode("\n      "));

const element19 = document.createElement("dt");
element19.append(document.createTextNode("Solution"));
element2.append(element19);

const element20 = document.createElement("dd");
element20.append(document.createTextNode("Paris"));
element2.append(element20);
element2.append(document.createTextNode("\n    "));
element1.append(element2);
element1.append(document.createTextNode("\n  "));
example.append(element1);

document.body.append(example);
|
\end{document}

% !tp-example Derivative of tan(x)
\documentclass{article}
\begin{document}
\wcTpSingleChoiceQuestion[answer=2,random,lang=en]{
  \begin{description}
    \item[{Title}] Derivative of the tangent function
    \item[{Prompt}] What is the derivative of $f(x) = \tan(x)$ at a point where $\cos(x) \ne 0$?
    \item[{Form}] \begin{itemize}
      \item $\cos^{2}(x)$
      \item $\displaystyle \frac{1}{\cos^{2}(x)}$
      \item $\displaystyle -\frac{1}{\sin^{2}(x)}$
      \item $\displaystyle \frac{1}{\tan^{2}(x)}$
    \end{itemize}
    \item[{Feedback}] \begin{itemize}
      \item The square of the cosine belongs in the denominator.
      \item Correct. Apply the quotient rule to sine divided by cosine.
      \item This is the derivative of the cotangent, not the tangent.
      \item The reciprocal square of the tangent is not its derivative.
    \end{itemize}
    \item[{Solution}] $\displaystyle f'(x) = \frac{1}{\cos^{2}(x)} = 1 + \tan^{2}(x)$, wherever $\cos(x) \ne 0$.
  \end{description}
}

\wcScript|
  // Wait for the question and its radio controls before typesetting the live content.
  window.MathJax = {
    svg: { fontCache: "local" },
    startup: {
      pageReady: async () => {
        await customElements.whenDefined("tp-single-choice-question");
        await customElements.whenDefined("tp-radio-list");
        return MathJax.startup.defaultPageReady();
      }
    }
  };
|

\wcScript[defer,src=https://cdn.jsdelivr.net/npm/mathjax@4/tex-svg.js]||
\end{document}

% !tp-example Regular heptagon
\documentclass{article}
\begin{document}
\wcTpSingleChoiceQuestion[answer=3,random,lang=en]{
  \begin{description}
    \item[{Title}] Recognize a regular heptagon
    \item[{Prompt}] Which figure is a regular heptagon? Select the polygon with seven equal sides and seven equal interior angles.
    \item[{Form}] \begin{itemize}
      \item \wcImg*[src=/docs/components/single-choice-question/examples/pentagon.svg,alt=Regular polygon with 5 sides,width=120,height=120]{}
      \item \wcImg*[src=/docs/components/single-choice-question/examples/hexagon.svg,alt=Regular polygon with 6 sides,width=120,height=120]{}
      \item \wcImg*[src=/docs/components/single-choice-question/examples/heptagon.svg,alt=Regular polygon with 7 sides,width=120,height=120]{}
      \item \wcImg*[src=/docs/components/single-choice-question/examples/octagon.svg,alt=Regular polygon with 8 sides,width=120,height=120]{}
      \item \wcImg*[src=/docs/components/single-choice-question/examples/nonagon.svg,alt=Regular polygon with 9 sides,width=120,height=120]{}
    \end{itemize}
    \item[{Feedback}] \begin{itemize}
      \item This pentagon has five sides.
      \item This hexagon has six sides.
      \item Correct. This heptagon has seven equal sides and seven equal interior angles.
      \item This octagon has eight sides.
      \item This nonagon has nine sides.
    \end{itemize}
    \item[{Solution}] The regular heptagon is the polygon with seven equal sides and seven equal interior angles.
  \end{description}
}
\end{document}
